\chapter{Ruang \texorpdfstring{$L^p$}{Lp}}\label{ch:b3:lp}

Integral Lebesgue dibangun untuk analisis; dan ruang $L^p$
adalah tempat analisis itu tinggal. Ruang itu Banach
(menurut Riesz--Fischer) --- yakni \omterm{thm:b3:complete:completion}{pelengkapan} yang menurut
\cref{exo:b3:complete:1} tak dipunyai fungsi
\omterm{def:b3:topology:continuity}{kontinu} --- dan ia menopang teknologi pemulusan,
yakni \emph{konvolusi terhadap \omterm{def:b3:lp:mollifier}{pemulus}}, yang menghampiri setiap
fungsi $L^p$ oleh fungsi $\mathcal C^\infty$. Bab ini membuktikan
versi integral Hölder dan Minkowski, \omterm{def:b3:complete:complete}{kelengkapannya},
teorema \omterm{ex:b3:lebesgue:gamma}{kepadatannya}, dan mesin peregularannya, lalu berakhir
dengan geografi keterkandungan dan interpolasi tangga
$L^p$. Sepanjang bab, $(X, \mathcal A, \mu)$ adalah \omterm{def:b3:measure:measure}{ruang ukuran}
dan fungsinya bernilai kompleks; sedangkan pada $\R^d$, \omterm{def:b3:measure:measure}{ukurannya} adalah
$\lambda_d$.

\section{Definisi; Hölder dan Minkowski}

\begin{definition}\label{def:b3:lp:space}
Untuk $1 \leq p < \infty$, $\mathcal L^p(\mu)$ adalah himpunan
$f$ yang \omterm{def:b3:lebesgue:measurable}{terukur} dengan $\norm f_p = \bigl(\int\abs
f^p\dd\mu\bigr)^{1/p} < \infty$, sedangkan $\mathcal L^\infty(\mu)$
himpunan $f$ yang terbatas di luar sebuah himpunan nol, dengan $\norm
f_\infty$ berupa \emph{sup esensialnya} --- yakni $M$ terkecil dengan
$\abs f \leq M$ hampir di mana-mana (dan infimumnya tercapai: iriskan himpunan
nolnya untuk $M + \frac1n$). Karena $\norm f_p = 0$ hanya memaksa $f
= 0$ \emph{hampir di mana-mana} (\cref{exo:b3:lebesgue:5}), kita mendefinisikan
\[
L^p(\mu) = \mathcal L^p(\mu)/\{f = 0 \text{ hampir di mana-mana}\} :
\]
yakni unsurnya berupa kelas fungsi modulo himpunan nol, dan
$\norm\cdot_p$ merupakan norma sejati pada
$L^p$.\index{ruang Lp@ruang $L^p$}
\end{definition}

\begin{theorem}[Ketaksamaan Hölder]\label{thm:b3:lp:holder}
Misalkan $1 \leq p, q \leq \infty$ dengan $\frac1p + \frac1q = 1$
(yakni eksponen sekawan). Untuk $f, g$ yang \omterm{def:b3:lebesgue:measurable}{terukur}:\index{ketaksamaan
Hölder}
\[
\norm{fg}_1 \leq \norm f_p\,\norm g_q ,
\]
dengan kesamaannya (untuk $1 < p < \infty$, norma berhingga, dan $f,g\ne0$)
jika dan hanya jika $\abs f^p$ dan $\abs g^q$ sebanding hampir di mana-mana.
\end{theorem}

\begin{proof}
Kasus $\{p, q\} = \{1, \infty\}$ berlangsung langsung (sebab $\abs{fg} \leq
\norm g_\infty\abs f$ hampir di mana-mana). Misalkan $1 < p < \infty$; lalu normalkan
$\norm f_p = \norm g_q = 1$ (lewat kehomogenannya; sedangkan norma nol atau tak
berhingga bersifat trivial). Ketaksamaan Young $ab \leq \frac{a^p}p +
\frac{b^q}q$ (dengan $a, b \geq 0$; lewat kecekungan $\ln$, seperti pada
\cref{pb:b3:banach:1}) memberikan secara titik demi titik $\abs{f g} \leq
\frac{\abs f^p}p + \frac{\abs g^q}q$; lalu integralkan: $\norm{fg}_1
\leq \frac1p + \frac1q = 1$. Kesamaannya memaksa kesamaan hampir di mana-mana pada
Young, yakni $\abs f^p = \abs g^q$ hampir di mana-mana (setelah
penormalannya; dan setelah dibatalkan, kesebandingannya).
\end{proof}

\begin{theorem}[Ketaksamaan Minkowski]\label{thm:b3:lp:minkowski}
Untuk $1 \leq p \leq \infty$: $\norm{f + g}_p \leq \norm f_p +
\norm g_p$.\index{ketaksamaan Minkowski}
\end{theorem}

\begin{proof}
Untuk $p = 1, \infty$: lewat ketaksamaan segitiga titik demi titik atau hampir di mana-mana. Sedangkan untuk $1 <
p < \infty$, andaikan $\norm{f+g}_p < \infty$ (jika tidak, pakailah
$\abs{f+g}^p \leq 2^{p-1}(\abs f^p + \abs g^p)$, dari kecembungan
$t^p$, untuk melihat bahwa ruas kirinya berhingga bila ruas kanannya demikian).
Maka
\[
\norm{f{+}g}_p^p \leq \int\abs f\,\abs{f{+}g}^{p-1} +
\int\abs g\,\abs{f{+}g}^{p-1}
\leq \bigl(\norm f_p + \norm g_p\bigr)\,
\bigl\|\abs{f{+}g}^{p-1}\bigr\|_q
\]
lewat Hölder, dan $\norm{\abs{f+g}^{p-1}}_q =
\norm{f+g}_p^{p/q}$ sebab $(p-1)q = p$; lalu bagilah dengan
$\norm{f+g}_p^{p/q}$ (jika tak nol; jika nol maka trivial) lalu pakai $p -
\frac pq = 1$.
\end{proof}

\section{Kelengkapan dan kawan-kawannya}

\begin{theorem}[Riesz--Fischer]\label{thm:b3:lp:complete}
Untuk $1 \leq p \leq \infty$, $L^p(\mu)$ merupakan ruang Banach.
Lebih lanjut, setiap barisan yang konvergen di $L^p$ mempunyai subbarisan
yang konvergen \emph{hampir di mana-mana} (dengan pendominasi $L^p$
pada kasus $p < \infty$).\index{teorema Riesz--Fischer}
\end{theorem}

\begin{proof}
Untuk $p = \infty$: barisan Cauchy terhadap $\norm\cdot_\infty$ bersifat, di luar
satu himpunan nol (yakni gabungan terbilang banyak), Cauchy seragam:
sehingga ia konvergen seragam di luarnya; selesai. Misalkan $p < \infty$. Menurut
\cref{exo:b3:complete:1}(b) cukup dijumlahkan deret yang konvergen
mutlak: misalkan $\sum\norm{f_k}_p = M < \infty$. Tetapkan
$G_n = \sum_{k \leq n}\abs{f_k}$ dan $G = \sum_k\abs{f_k}$
(secara titik demi titik di $[0,\infty]$): maka menurut Minkowski $\norm{G_n}_p \leq
M$, dan kekonvergenan monoton (sebab $G_n^p \nearrow G^p$) memberikan $\int G^p \leq M^p$:
sehingga $G < \infty$ hampir di mana-mana, jadi deret $\sum f_k(x)$ konvergen
mutlak untuk hampir setiap $x$; sebutlah jumlahnya $S(x)$ (dengan nilai apa pun pada
himpunan nolnya). Maka $\abs{S - \sum_{k\leq n}f_k}^p \leq (2G)^p
\in L^1$, dan kekonvergenan terdominasinya memberikan $\norm{S - \sum_{k\leq n}f_k}_p \to 0$:
jadi deretnya konvergen di $L^p$.

Untuk pernyataan subbarisannya: jika $f_n \to f$ di $L^p$, pilihlah $n_k$
dengan $\norm{f_{n_{k+1}} - f_{n_k}}_p \leq 2^{-k}$; maka deret
$\sum(f_{n_{k+1}} - f_{n_k})$ jatuh di bawah argumen
sebelumnya: ia konvergen mutlak hampir di mana-mana, terdominasi oleh sebuah $G \in
L^p$, sehingga $f_{n_k} \to f_{n_1} + \sum(\cdots)$ hampir di mana-mana, dan limit
hampir di mana-mana itu haruslah (wakil dari) $f$ (sebab keduanya limit
$L^p$). Pendominasinya: $\abs{f_{n_k}} \leq \abs{f_{n_1}} + G$.
\end{proof}

\begin{remark}\label{rem:b3:lp:modes}
Kekonvergenan $L^p$ tak mengakibatkan kekonvergenan hampir di mana-mana (lihat
barisan \emph{mesin tik}, \cref{exo:b3:lp:3}), dan sebaliknya pun tidak
(lihat gundukan yang lolos): sebab kedua ragamnya terkait hanya
lewat subbarisan dan dominasi. Menyimpan contoh penyangkal
\cref{exo:b3:lp:3} dalam ingatan adalah vaksin yang terbaik.
\end{remark}

\section{Teorema kepadatan}

\begin{theorem}\label{thm:b3:lp:density}
Misalkan $1 \leq p < \infty$.
\begin{enumerate}
  \item \omterm{def:b3:lebesgue:simple}{Fungsi sederhana} (dengan pendukung berukuran berhingga) bersifat
        padat di $L^p(\mu)$.
  \item Di $L^p(\R^d)$, fungsi \omterm{def:b3:topology:continuity}{kontinu} berpendukung \omterm{def:b3:topology:compact}{kompak}
        $\mathcal C_c(\R^d)$ bersifat padat.
  \item Translasi bersifat \omterm{def:b3:topology:continuity}{kontinu} pada $L^p(\R^d)$: dengan menulis
        $\tau_hf = f(\cdot - h)$, berlaku $\norm{\tau_hf - f}_p \to 0$
        saat $h \to 0$.
\end{enumerate}
Tak satu pun dari ketiganya berlaku untuk $p = \infty$.
\end{theorem}

\begin{proof}
(1) Untuk $f \geq 0$: barisan diadik $s_n \nearrow f$ pada
\cref{thm:b3:lebesgue:approximation} memenuhi $\abs{f - s_n}^p
\leq f^p \in L^1$: lewat kekonvergenan terdominasi. (Setiap $s_n \leq f$ terletak di $L^p$, dan
himpunan arasnya berukuran berhingga di tempat nilainya
positif: sebab $\mu(s_n \geq c) \leq c^{-p}\int f^p$.) Lalu pecahlah
$f$ yang umum menjadi empat bagian tak negatif.

(2) Menurut (1) cukup dihampiri $\mathbf 1_A$ dengan $A$ Borel
yang $\lambda_d(A) < \infty$. Keteraturannya (dengan bukti seperti pada
\cref{thm:b3:measure:regularity}) memberikan $K \subseteq A
\subseteq U$, \omterm{def:b3:topology:compact}{kompak} di dalam terbuka, dengan $\lambda_d(U\setminus K) <
\varepsilon$; lalu fungsi Urysohn
\[
\varphi(x) = \frac{d(x, \R^d\setminus U)}{d(x, \R^d\setminus U)
+ d(x, K)}
\]
bersifat \omterm{def:b3:topology:continuity}{kontinu}, bernilai $1$ pada $K$ dan $0$ di luar $U$, dan dapat diambil
berpendukung \omterm{def:b3:topology:compact}{kompak} (susutkanlah $U$ menjadi \omterm{def:b3:topology:topology}{himpunan terbuka} terbatas lebih dahulu). Maka
$\norm{\mathbf 1_A - \varphi}_p^p \leq \lambda_d(U\setminus K)
< \varepsilon$.

(3) Untuk $g \in \mathcal C_c$: \omterm{def:b3:topology:continuity}{kekontinuan} seragamnya memberikan
$\norm{\tau_hg - g}_\infty \to 0$, dengan pendukungnya di dalam satu
\omterm{def:b3:topology:compact}{kompak} tetap bila $\abs h \leq 1$: sehingga $\norm{\tau_hg - g}_p \to 0$. Sedangkan untuk
$f$ yang umum: pilihlah $g \in \mathcal C_c$ dengan $\norm{f - g}_p <
\varepsilon$; maka $\norm{\tau_hf - f}_p \leq 2\norm{f - g}_p +
\norm{\tau_hg - g}_p$ (lewat keinvarianan translasi normanya).

Untuk $p = \infty$: penghampiran seragam atas $\mathbf
1_{\intoo0\infty}$ oleh fungsi \omterm{def:b3:topology:continuity}{kontinu} mustahil
(sebab ada lompatan), dan $\norm{\tau_h\mathbf 1_{\intoo0\infty} - \mathbf
1_{\intoo0\infty}}_\infty = 1$ untuk $h \neq 0$.
\end{proof}

\section{Konvolusi dan peregularan}

\begin{theorem}[Ketaksamaan Young]\label{thm:b3:lp:young}
Misalkan $1 \leq p \leq \infty$, $f \in L^1(\R^d)$, dan $g \in
L^p(\R^d)$. Maka $f * g$ terdefinisi hampir di mana-mana, termasuk $L^p$,
dan\index{ketaksamaan Young}
\[
\norm{f * g}_p \leq \norm f_1\,\norm g_p .
\]
\end{theorem}

\begin{proof}
Untuk $p = \infty$: lewat batas langsungnya. Untuk $p = 1$:
\cref{thm:b3:product:convolution}. Misalkan $1 < p < \infty$ dengan $q$
sekawannya. Pecahlah $\abs{f(y)} = \abs{f(y)}^{1/q}\cdot
\abs{f(y)}^{1/p}$ lalu terapkan Hölder:
\[
\int\abs{f(y)}\,\abs{g(x{-}y)}\,\dd y
\leq \Bigl(\int\abs f\Bigr)^{1/q}
\Bigl(\int\abs{f(y)}\,\abs{g(x - y)}^p\,\dd y\Bigr)^{1/p} .
\]
Lalu pangkatkan ke $p$ dan integralkan terhadap $x$; maka Tonelli pada
faktor keduanya memberikan $\norm f_1^{p/q}\cdot\norm f_1\norm g_p^p$,
yakni $\norm{f*g}_p^p \leq \norm f_1^{1 + p/q}\norm g_p^p =
(\norm f_1\norm g_p)^p$ --- sedangkan keberhinggaan integral
Tonellinya membenarkan kekonvergenan mutlak hampir di mana-mana seperti pada
\cref{thm:b3:product:convolution}.
\end{proof}

\begin{definition}[Pemulus]\label{def:b3:lp:mollifier}
Fungsi\index{pemulus}\index{fungsi gundukan}
\[
\rho(x) = \begin{cases}
c\,\exp\Bigl(-\dfrac{1}{1 - \norm x^2}\Bigr) & \norm x < 1,\\
0 & \norm x \geq 1,
\end{cases}
\]
dengan $c$ yang menormalkan $\int\rho = 1$, bersifat $\mathcal C^\infty$ pada
$\R^d$: sebab intinya adalah bahwa $t \mapsto \eu^{-1/t}\mathbf
1_{t>0}$ bersifat $\mathcal C^\infty$ pada $\R$, dengan semua turunannya di
$0^+$ bernilai $0$ (sebab setiap turunannya berbentuk $P(1/t)\eu^{-1/t}$ untuk sebuah
polinomial $P$, yang menuju $0$; lewat induksi). Untuk
$\varepsilon > 0$ tetapkan $\rho_\varepsilon(x) =
\varepsilon^{-d}\rho(x/\varepsilon)$: yang berpendukung di $\bar B(0,
\varepsilon)$, dan tetap berintegral $1$.
\end{definition}

\begin{theorem}[Peregularan]\label{thm:b3:lp:regularization}
Misalkan $1 \leq p < \infty$ dan $f \in L^p(\R^d)$. Maka:
\begin{enumerate}
  \item $f * \rho_\varepsilon \in \mathcal C^\infty(\R^d)$,
        dengan $\partial^\alpha(f * \rho_\varepsilon) = f *
        \partial^\alpha\rho_\varepsilon$;
  \item $\norm{f * \rho_\varepsilon - f}_p \to 0$ saat
        $\varepsilon \to 0$;
  \item sehingga $\mathcal C^\infty_c(\R^d)$ bersifat padat di
        $L^p(\R^d)$.
\end{enumerate}
\end{theorem}

\begin{proof}
(1) Lewat pendiferensialan di bawah integralnya
(\cref{thm:b3:lebesgue:paramdiff}) terhadap $x$: sebab untuk $x$ di sebuah bola
$B$, $\abs{\partial_{x_i}\rho_\varepsilon(x - y)} \leq
C_\varepsilon\,\mathbf 1_{K}(y)$ dengan $K$ \omterm{def:b3:topology:compact}{kompak} (yakni $y$ sejauh
$\varepsilon$ dari $B$), dan $\abs f\,\mathbf 1_K \in L^1$
(lewat Hölder terhadap $\mathbf 1_K$): sehingga teoremanya berlaku;
lalu iterasikan untuk turunan yang lebih tinggi.

(2) Karena $\int\rho_\varepsilon = 1$:
\[
(f * \rho_\varepsilon)(x) - f(x)
= \int \bigl(f(x - y) - f(x)\bigr)\rho_\varepsilon(y)\,\dd y ,
\]
dan ketaksamaan \emph{integral} Minkowski --- atau langsung:
Hölder/Jensen dengan \omterm{def:b3:measure:measure}{ukuran} peluang
$\rho_\varepsilon\dd y$ dan Tonelli ---
\[
\norm{f*\rho_\varepsilon - f}_p^p
\leq \int\Bigl(\int\abs{f(x-y) -
f(x)}^p\dd x\Bigr)\rho_\varepsilon(y)\,\dd y
= \int \norm{\tau_yf - f}_p^p\;\rho_\varepsilon(y)\,\dd y
\]
(dengan langkah tengahnya: terapkan ketaksamaan Jensen,
\cref{exo:b3:lp:10}, pada integral-$y$ dalamnya, lalu Tonelli).
Integrannya berpendukung di $\norm y \leq \varepsilon$ dan
menuju $0$ di sana secara seragam saat $\varepsilon \to 0$
(\cref{thm:b3:lp:density}(3)): sehingga seluruh ungkapannya menuju
$0$.

(3) Hampirilah $f$ oleh $g \in \mathcal C_c$
(\cref{thm:b3:lp:density}(2)), lalu $g$ oleh $g *
\rho_\varepsilon \in \mathcal C_c^\infty$ (dengan pendukung \omterm{def:b3:topology:compact}{kompak}:
yakni jumlah pendukungnya).
\end{proof}

\begin{example}[Memuluskan $\abs x$, beserta lajunya]
\label{ex:b3:lp:mollifyexample}
Ambillah $f(x) = \abs x$ pada $\R$ (yang $L^1$ secara lokal; sehingga teoremanya
berlaku pada setiap jendela terbatas) dan sebuah \omterm{def:b3:lp:mollifier}{pemulus} setangkup
$\rho_\varepsilon$. Maka
\[
f_\varepsilon(x) = (f * \rho_\varepsilon)(x)
= \int\abs{x - y}\,\rho_\varepsilon(y)\,\dd y
\]
bersifat $\mathcal C^\infty$; sedangkan jauh dari tekukannya, tak terjadi apa-apa:
sebab untuk $\abs x \geq \varepsilon$, $\abs{x - y}$ linear terhadap $x$
pada pendukung $\rho_\varepsilon$, sehingga $f_\varepsilon(x) =
\abs x$ \emph{secara persis} (sebab kesetangkupannya membunuh koreksinya).
Di dekat $0$, pemulusannya berbiaya tepat
\[
0 \leq f_\varepsilon(0) = \int\abs
y\,\rho_\varepsilon(y)\,\dd y \leq \varepsilon,
\qquad
\norm{f_\varepsilon - f}_\infty \leq \varepsilon :
\]
sehingga galat penghampirannya terkurung di
persekitaran-$\varepsilon$ kesingularannya dan seukuran
dengannya. Sementara itu $f_\varepsilon'' \geq 0$ di mana-mana
(sebab $f$ cembung, dan konvolusi terhadap $\rho_\varepsilon
\geq 0$ mengawetkan kecembungan), dengan $\int f_\varepsilon'' =
f_\varepsilon'(\infty) - f_\varepsilon'(-\infty) = 2$: jadi
turunan keduanya berupa gundukan bermassa $2$ yang terjepit ke dalam lebar
$O(\varepsilon)$, sehingga $\norm{f_\varepsilon''}_\infty \gtrsim
\varepsilon^{-1}$. Jadi pemulusan adalah tukar tambah: galat seragam
$O(\varepsilon)$ ditukar dengan peledakan turunan
$O(\varepsilon^{-1})$ --- yakni kurs tukar persis yang
diformalkan analisis kuantitatif (lewat ketaksamaan interpolasi,
yakni lingkaran gagasan \cref{pb:b3:lp:1}).
\end{example}

\begin{corollary}[Lema fundamental kalkulus
variasi]\label{cor:b3:lp:fundlemma}
Misalkan $f \in L^1_{\mathrm{loc}}(\R^d)$ (yakni \omterm{def:b3:lebesgue:l1}{terintegralkan} pada \omterm{def:b3:topology:compact}{kompak})
dengan $\int f\varphi = 0$ untuk setiap $\varphi \in \mathcal
C^\infty_c(\R^d)$. Maka $f = 0$ hampir di mana-mana.\index{lema fundamental
kalkulus variasi}
\end{corollary}

\begin{proof}
Tetapkan sebuah bola $B = B(0, R)$ lalu misalkan $g = f\mathbf 1_{B(0, R+1)}
\in L^1$. Untuk $x \in B$ dan $\varepsilon < 1$: $(g *
\rho_\varepsilon)(x) = \int f(y)\rho_\varepsilon(x - y)\dd y =
0$, sebab fungsi ujinya adalah $y \mapsto \rho_\varepsilon(x-y)
\in \mathcal C_c^\infty$. Padahal $g * \rho_\varepsilon \to g$ di
$L^1$ (\cref{thm:b3:lp:regularization}): sehingga $g = 0$ hampir di mana-mana pada $B$;
lalu habiskanlah $\R^d$.
\end{proof}

\section{Geografi \texorpdfstring{$L^p$}{Lp}}

\begin{proposition}\label{prop:b3:lp:inclusions}
(a) Jika $\mu(X) < \infty$ dan $1 \leq p \leq q \leq \infty$,
maka $L^q \subseteq L^p$ dengan $\norm f_p \leq
\mu(X)^{\frac1p - \frac1q}\,\norm f_q$.
(b) Pada $\R^d$ (yang berukuran tak berhingga) tak ada keterkandungan:
sebab untuk $p \neq q$ ada fungsi di $L^p\setminus L^q$.
(c) (Interpolasi) Jika $p < r < q$ dan $\alpha \in \intoo01$
didefinisikan oleh $\frac1r = \frac\alpha p + \frac{1 - \alpha}q$,
maka
\[
\norm f_r \leq \norm f_p^{\alpha}\,\norm f_q^{1 - \alpha} ;
\]
khususnya $L^p \cap L^q \subseteq L^r$.
\end{proposition}

\begin{proof}
(a) Lewat Hölder dengan eksponen $\frac qp$ dan sekawannya:
$\int\abs f^p\cdot 1 \leq \norm{\abs f^p}_{q/p}\,\norm
1_{(q/p)'} = \norm f_q^p\,\mu(X)^{1 - p/q}$ (sedangkan $q = \infty$
langsung). (b) Di dekat $0$ dan di dekat $\infty$, pangkat $x^{-\alpha}$
mengkalibrasinya: \cref{exo:b3:lp:2}. (c) Tulislah $\abs f^r = \abs
f^{r\alpha}\abs f^{r(1-\alpha)}$ lalu terapkan Hölder dengan
pasangan sekawan $\frac p{r\alpha}$ dan $\frac q{r(1-\alpha)}$
(yang sekawan persis menurut definisi $\alpha$):
$\int\abs f^r \leq \norm f_p^{r\alpha}\norm f_q^{r(1 -
\alpha)}$.
\end{proof}

\begin{method}\label{met:b3:lp:toolkit}
Perkakas $L^p$, sebagaimana dipakai di mana-mana di bawah ini: untuk membuktikan sebuah
identitas atau ketaksamaan bagi setiap $f \in L^p$ --- buktikanlah pada sebuah
kelas yang padat ($\mathcal C_c^\infty$ lewat
\cref{thm:b3:lp:regularization}) lalu perluas lewat \omterm{def:b3:topology:continuity}{kekontinuannya}
(\cref{thm:b3:complete:extension}, sebab kedua ruasnya
kontinu-$L^p$); untuk membuktikan $f = 0$, ujilah terhadap $\mathcal
C_c^\infty$ (\cref{cor:b3:lp:fundlemma}); untuk memperoleh kemulusan,
konvolusikanlah; dan untuk menukar eksponennya, pakailah Hölder dan interpolasi. Sedangkan
teori Fourier \cref{ch:b3:fouriertransform} merupakan satu terapan panjang
metode ini.
\end{method}

\section{Latihan}

\begin{exercise}[$\star$]\label{exo:b3:lp:1}
(a) Nyatakan dan buktikan ketaksamaan Cauchy--Schwarz di
$L^2(\mu)$ sebagai kasus $p = q = 2$ dari Hölder.
(b) Pada ruang \emph{peluang}, tunjukkan bahwa $p \mapsto \norm f_p$
bersifat tak turun.
(c) Kapankah Hölder menjadi kesamaan untuk $p = 1$ dan $q = \infty$?
\end{exercise}

\begin{exercise}[$\star$]\label{exo:b3:lp:2}
Untuk $p \in \intco1\infty$ yang mana fungsi berikut termasuk
$L^p$?
\[
x^{-1/2}\mathbf 1_{\intoo01},\qquad
x^{-1/2}\mathbf 1_{\intoo1\infty},\qquad
\frac{1}{x^{1/2}(1 + \abs{\ln x})}\ \text{pada } \intoo01,
\qquad
\frac1{1 + \abs x}\ \text{pada } \R .
\]
Simpulkan: pada $\intoo01$ $p$ yang kecil lebih mudah, pada
$\intoo1\infty$ $p$ yang besar lebih mudah, dan tak ada $L^p$ yang memuat
yang lain pada $\R$.
\end{exercise}

\begin{exercise}[$\star\star$]\label{exo:b3:lp:3}
(Mesin tik) Cacahlah selang diadiknya $I_1 =
\intcc01$, $I_2 = \intcc0{\frac12}$, $I_3 = \intcc{\frac12}1$,
$I_4 = \intcc0{\frac14}$, \dots\ lalu misalkan $f_n = \mathbf
1_{I_n}$.
(a) Tunjukkan bahwa $f_n \to 0$ di setiap $L^p(\intcc01)$ dengan $p < \infty$,
tetapi $(f_n(x))$ divergen untuk \emph{setiap} $x \in \intcc01$.
(b) Tampilkan subbarisan yang konvergen hampir di mana-mana yang dijanjikan
\cref{thm:b3:lp:complete}.
(c) Sebaliknya berikan sebuah barisan yang konvergen hampir di mana-mana tetapi tak di
$L^1$, dan satu yang konvergen di $L^1$ tetapi tak di $L^p$ mana pun dengan $p > 1$.
\end{exercise}

\begin{exercise}[$\star\star$]\label{exo:b3:lp:4}
Misalkan $\mu(X) < \infty$ dan $f \in L^\infty(\mu)$ dengan $f \neq 0$.
Tunjukkan bahwa $\norm f_p \to \norm f_\infty$ saat $p \to \infty$.
\emph{(Batas atasnya lewat (a) dari \cref{prop:b3:lp:inclusions};
sedangkan batas bawahnya dengan mengintegralkan atas $\{\abs f > \norm f_\infty -
\varepsilon\}$, yang berukuran positif.)}
\end{exercise}

\begin{exercise}[$\star\star$]\label{exo:b3:lp:5}
(a) Di manakah tepatnya bukti
\cref{thm:b3:lp:density}(3) memakai $p < \infty$?
(b) Tunjukkan bahwa $f \in L^\infty(\R)$ memenuhi $\norm{\tau_hf -
f}_\infty \to 0$ jika dan hanya jika $f$ mempunyai wakil yang \omterm{def:b3:topology:continuity}{kontinu}
seragam.
\end{exercise}

\begin{exercise}[$\star\star$]\label{exo:b3:lp:6}
Misalkan $p, q$ sekawan, $f \in L^p(\R^d)$, dan $g \in
L^q(\R^d)$. Tunjukkan bahwa $f * g$ terdefinisi \emph{di mana-mana},
terbatas, dengan $\norm{f*g}_\infty \leq \norm f_p\norm g_q$, dan
\emph{\omterm{def:b3:topology:continuity}{kontinu} seragam}. \emph{(Pakai \omterm{def:b3:topology:continuity}{kekontinuan} translasi
di $L^p$; lalu tanganilah $p \in \{1, \infty\}$ secara terpisah --- sebab untuk $p =
\infty$ pakailah \omterm{def:b3:topology:continuity}{kekontinuan} translasi pada faktor $L^1$-nya.)}
\end{exercise}

\begin{exercise}[$\star\star$]\label{exo:b3:lp:7}
Misalkan $f \in L^1_{\mathrm{loc}}(\intoo ab)$ dengan
$\int f\varphi' = 0$ untuk setiap $\varphi \in \mathcal
C^\infty_c(\intoo ab)$. Tunjukkan bahwa $f$ sama dengan sebuah konstanta
hampir di mana-mana. \emph{(Tetapkan $\chi \in \mathcal C_c^\infty$ dengan
$\int\chi = 1$; maka setiap $\psi \in \mathcal C_c^\infty$ dengan
$\int\psi = 0$ berupa sebuah $\varphi'$; lalu tulislah fungsi uji umumnya
sebagai $\psi + (\int\psi)\chi$ lalu terapkan
\cref{cor:b3:lp:fundlemma} pada $f - c$ dengan $c = \int
f\chi$.)}
\end{exercise}

\begin{exercise}[$\star\star\star$]\label{exo:b3:lp:8}
(Urysohn yang mulus) Misalkan $K \subseteq U \subseteq \R^d$ dengan $K$
\omterm{def:b3:topology:compact}{kompak} dan $U$ terbuka. Konstruksikan $\varphi \in \mathcal
C^\infty_c(\R^d)$ dengan $0 \leq \varphi \leq 1$, $\varphi = 1$
pada $K$, dan $\operatorname{supp}\varphi \subseteq U$.
\emph{(Muluskanlah indikator persekitaran-$\delta$
$K_\delta$ dari $K$ dengan $\rho_{\delta/2}$, untuk $\delta$ yang kecil.)}
Lalu turunkan pernyataan partisi satuan $\mathcal C^\infty$ bagi
sebuah \omterm{def:b3:topology:compact}{kompak} yang terselimuti berhingga banyak \omterm{def:b3:topology:topology}{himpunan terbuka}.
\end{exercise}

\begin{exercise}[$\star\star$]\label{exo:b3:lp:9}
Dengan memakai interpolasi (\cref{prop:b3:lp:inclusions}(c)):
(a) tunjukkan bahwa $L^1(\R)\cap L^\infty(\R) \subseteq L^p(\R)$ untuk
setiap $p$, dengan $\norm f_p \leq \norm f_1^{1/p}\norm
f_\infty^{1 - 1/p}$;
(b) tunjukkan bahwa $f \mapsto \norm f_p$ bersifat, untuk $f$ yang tetap,
log-cembung terhadap $\frac1p$, lalu berikan contoh yang $f \in L^p$-nya
tepat untuk $p$ pada suatu selang $(p_0, p_1)$.
\end{exercise}

\begin{exercise}[$\star\star$]\label{exo:b3:lp:10}
(Jensen) Misalkan $\mu$ sebuah \omterm{def:b3:measure:measure}{ukuran} \emph{peluang}, $f \in
L^1(\mu)$ real, dan $\Phi \colon \R \to \R$ cembung. Tunjukkan bahwa
\[
\Phi\Bigl(\int f\,\dd\mu\Bigr) \leq \int \Phi\circ
f\,\dd\mu
\]
\emph{(pakailah garis penopang $\Phi$ di titik $m = \int f$)}.
Lalu turunkan ketaksamaan rata-rata hitung--ukur beserta
kemonotonan $p \mapsto \norm f_p$ pada
\cref{exo:b3:lp:1}(b).\index{ketaksamaan Jensen}
\end{exercise}

\begin{exercise}[$\star\star\star$]\label{exo:b3:lp:11}
(Ketaksamaan konvolusi Young) Misalkan $1 \leq p, q, r \leq
\infty$ dengan $\frac1p + \frac1q = 1 + \frac1r$, dan $f \in
L^p(\R^d)$, $g \in L^q(\R^d)$.
(a) Buktikan $\norm{f * g}_r \leq \norm f_p\,\norm g_q$.
\emph{(Tulislah, untuk eksponen sekawan yang dihitung dari $p, q,
r$,
\[
\abs{f(y)g(x-y)} = \bigl(\abs f^p\abs g^q\bigr)^{1/r}
\cdot\abs f^{\,p(1/p - 1/r)}\cdot\abs g^{\,q(1/q - 1/r)},
\]
lalu terapkan ketaksamaan Hölder tiga faktor dengan
eksponen $r$, $\frac{pr}{r - p}$, $\frac{qr}{r - q}$; lalu
integralkan terhadap $x$ lewat Tonelli.)}
(b) Periksalah ketiga kasus khusus yang sudah dikenal: $r =
\infty$ (yakni Hölder, \cref{exo:b3:lp:6}); $q = 1$
(yakni kestabilan-$L^p$ konvolusi oleh kernel \omterm{def:b3:lebesgue:l1}{terintegralkan});
dan $p = q = 1$ (yakni $L^1$ sebagai aljabar konvolusi,
\cref{thm:b3:product:convolution}).
(c) Mengapa tak ada ketaksamaan dengan $\frac1p + \frac1q < 1 +
\frac1r$? \emph{(Ujilah pada pendilatan $f_\lambda(x) =
f(\lambda x)$ lalu bandingkan penskalaan kedua ruasnya.)}
\end{exercise}

\begin{exercise}[$\star\star$]\label{exo:b3:lp:12}
(Kasus kesamaan) (a) Pada ketaksamaan Hölder
$\int\abs{fg} \leq \norm f_p\norm g_q$ (dengan $1 < p < \infty$),
tunjukkan bahwa kesamaannya berlaku jika dan hanya jika $\abs f^p$ dan $\abs g^q$
sebanding hampir di mana-mana. \emph{(Lacaklah kasus kesamaan ketaksamaan
Young $ab \leq \frac{a^p}p + \frac{b^q}q$, yang berupa $a^p
= b^q$.)}
(b) Pada ketaksamaan Minkowski $\norm{f + g}_p \leq \norm f_p
+ \norm g_p$ (dengan $1 < p < \infty$), tunjukkan bahwa kesamaannya dengan $f,
g \neq 0$ memaksa $g = cf$ hampir di mana-mana dengan $c > 0$.
(c) Bandingkan dengan $p = 1$ dan $p = \infty$: uraikanlah
kasus kesamaannya di sana (yang jauh lebih besar), lewat contoh.
\end{exercise}

\section{Soal: ketaksamaan Hardy}

\begin{problem}[{Soal akhir pekan --- ketaksamaan Hardy dan
konstanta tajamnya}]\label{pb:b3:lp:1}
Untuk $f \in L^p(\intoo0{+\infty})$ dengan $1 < p < \infty$, definisikan
\emph{operator Hardy}
\[
(Hf)(x) = \frac1x\int_0^x f(t)\,\dd t .
\]
Ketaksamaan Hardy (1920) menyatakan\index{ketaksamaan Hardy}
\[
\norm{Hf}_p \;\leq\; \frac{p}{p-1}\,\norm f_p ,
\]
dan konstanta $\frac p{p-1}$ bersifat optimal serta tak tercapai.
Soal ini membuktikan segalanya, lalu memperluasnya ke deret.

\medskip
\noindent\textbf{Bagian I --- Ketaksamaannya.} Andaikan dahulu $f
\geq 0$ \omterm{def:b3:topology:continuity}{kontinu} dengan pendukung \omterm{def:b3:topology:compact}{kompak} di
$\intoo0{+\infty}$, lalu misalkan $F(x) = \int_0^xf$.
\begin{enumerate}
  \item Tunjukkan bahwa $Hf \in L^p$: sebab di dekat $0$, $F$ lenyap pada sebuah
        \omterm{def:b3:topology:topology}{persekitaran} $0$; sedangkan di dekat $\infty$, $F$ terbatas, sehingga
        $(Hf)(x) = O(1/x)$, dan $x \mapsto \frac1x$ termasuk
        $L^p(\intoo1{+\infty})$ untuk $p > 1$.
  \item Integralkan secara parsial untuk menunjukkan bahwa
        \[
        \int_0^\infty \Bigl(\frac Fx\Bigr)^p\dd x
        = \frac{p}{p-1}\int_0^\infty\Bigl(\frac
        Fx\Bigr)^{p-1}f(x)\,\dd x .
        \]
        \emph{(Turunkan $x^{1-p}F^p$; sedangkan suku batasnya
        lenyap --- benarkanlah kedua ujungnya.)}
  \item Terapkan Hölder pada ruas kanannya lalu turunkan
        $\norm{Hf}_p \leq \frac p{p-1}\norm f_p$ untuk $f$ semacam itu.
  \item Perluaslah ke seluruh $L^p$: untuk $f \geq 0$, konstruksikan
        $f_n$ yang \omterm{def:b3:topology:continuity}{kontinu} dengan pendukung \omterm{def:b3:topology:compact}{kompak} di
        $\intoo0{+\infty}$ dan $0 \leq f_n \nearrow f$ hampir di mana-mana
        \emph{(pancunglah, lalu hampirilah secara monoton ---
        benarkan konstruksinya)}; maka $Hf_n \nearrow Hf$
        secara titik demi titik (lewat kekonvergenan monoton di dalam rata-ratanya) dan kekonvergenan monotonnya mengalihkan
        ketaksamaannya ke limitnya. Sedangkan untuk $f$ yang bertanda atau kompleks,
        simpulkan dengan $\abs{Hf} \leq H\abs f$.
\end{enumerate}

\medskip
\noindent\textbf{Bagian II --- Keoptimalannya.}
\begin{enumerate}[resume]
  \item Untuk $A > 1$ misalkan $f_A(t) = t^{-1/p}\,\mathbf
        1_{\intcc1A}(t)$. Hitunglah $\norm{f_A}_p^p = \ln A$
        lalu, untuk $1 \leq x \leq A$,
        \[
        (Hf_A)(x) = \frac p{p-1}\;x^{-1/p}\,
        \bigl(1 - x^{-(1 - 1/p)}\bigr).
        \]
  \item Turunkan $\liminf_{A\to\infty} \norm{Hf_A}_p/
        \norm{f_A}_p \geq \frac{p}{p-1}$, lalu simpulkan bahwa
        konstantanya optimal.
  \item Tunjukkan bahwa kesamaan $\norm{Hf}_p = \frac
        p{p-1}\norm f_p$ dengan $f \neq 0$ mustahil.
        \emph{(Lacaklah kasus kesamaan Hölder pada pertanyaan
        3: ia akan memaksa perilaku bertipe $f = cx^{-1/p}$, yang
        tak berada di $L^p$.)}
\end{enumerate}

\medskip
\noindent\textbf{Bagian III --- Ketaksamaan diskretnya.}
\begin{enumerate}[resume]
  \item Untuk $g \geq 0$ yang tak naik pada $(0,\infty)$ dan
        $a_n = g(n)$, bandingkanlah $\sum a_n^p$ dengan $\int g^p$,
        beserta rata-rata-$H$-nya, untuk menurunkan dari Bagian I
        \emph{ketaksamaan Hardy diskret}: yakni untuk $a_n \geq 0$,
        \[
        \sum_{n\geq1}\Bigl(\frac{a_1 + \dots +
        a_n}{n}\Bigr)^{p}
        \;\leq\;
        \Bigl(\frac{p}{p-1}\Bigr)^{p}\,\sum_{n\geq1}a_n^p
        \]
        --- buktikanlah dahulu untuk $(a_n)$ yang tak naik lewat
        pembandingan di atas, lalu susutkan kasus umumnya ke
        kasus tak naiknya lewat penyusunan ulang \emph{(terimalah, dengan
        pembenaran satu baris, bahwa menyortir $(a_n)$ secara
        turun hanya dapat menaikkan ruas kirinya sambil
        membiarkan ruas kanannya tetap)}.
  \item Turunkan: jika $\sum a_n^p < \infty$ maka rata-rata Cesàro
        $(a_n)$ kembali berada di $\ell^p$ --- lalu berikan sebuah
        contoh ($p = 2$) yang $(a_n) \in \ell^2$-nya tetapi $a_n$-nya
        tak terjumlahkan, padahal Hardy tetap mengendalikan rata-ratanya.
\end{enumerate}

\medskip
\noindent\textbf{Bagian IV --- Epilog.}
\begin{enumerate}[resume]
  \item Tunjukkan bahwa ketaksamaan Hardy gagal untuk $p = 1$: dengan
        $f = \mathbf 1_{\intcc01}$, hitunglah $Hf$ lalu amati bahwa
        $Hf \notin L^1$. Di manakah buktinya patah?
\end{enumerate}

\medskip
\noindent\textbf{Bagian V --- Fungsi maksimalnya, dan
teorema pendiferensialan Lebesgue.} Hardy merata-ratakan dari
titik asal; sedangkan Hardy--Littlewood merata-ratakan \emph{di sekitar setiap titik}.
Untuk $f \in L^1(\R)$ definisikan
\[
Mf(x) = \sup_{r>0}\ \frac1{2r}\int_{x-r}^{x+r}\abs
f\,\dd\lambda .
\]
\begin{enumerate}[resume]
  \item (Vitali, versi berhingga) Misalkan $B_1, \dots, B_N$
        selang terbuka. Tunjukkan bahwa ada subkeluarga saling lepas
        $B_{i_1}, \dots, B_{i_k}$ dengan $\bigcup_jB_j
        \subseteq \bigcup_l3B_{i_l}$, dengan $3B$ menandai
        selang berpusat sama dan berpanjang tiga kali lipat
        \emph{(secara rakus: berulang kali pilihlah selang terpanjang
        yang saling lepas dengan yang sudah terpilih)}.
  \item (Tipe lemah $(1,1)$) Tunjukkan bahwa untuk setiap $t > 0$,
        \[
        \lambda\bigl(\{Mf > t\}\bigr) \;\leq\;
        \frac3t\,\norm f_1 :
        \]
        sebab setiap $x$ dengan $Mf(x) > t$ memiliki sebuah selang berpusat
        $B_x$ dengan $\int_{B_x}\abs f > t\,\lambda(B_x)$; lalu ambillah
        sebuah \omterm{def:b3:topology:compact}{kompak} $K \subseteq \{Mf > t\}$ (lewat keteraturan dalamnya),
        selimuti ia oleh berhingga banyak $B_x$, terapkan pertanyaan 11,
        lalu habiskan.
  \item Hitunglah $M\mathbf 1_{\intcc01}(x)$ untuk $x > 1$ lalu
        turunkan bahwa $Mf \notin L^1$ untuk setiap $f \neq 0$
        (sebab $Mf(x) \geq \frac c{\abs x}$ di tak hingga): sehingga pada $p =
        1$, ketaksamaan lemah pertanyaan 12 merupakan pernyataan terbaik
        yang mungkin.
  \item (Tipe kuat untuk $p > 1$) Untuk $f \in L^p$: pecahlah $f =
        f\,\mathbf 1_{\abs f > t/2} + f\,\mathbf 1_{\abs f
        \leq t/2}$, amati bahwa $Mf \leq M\bigl(f\mathbf 1_{\abs
        f > t/2}\bigr) + \frac t2$, lalu gabungkan pertanyaan 12
        dengan rumus kue berlapis
        (\cref{prop:b3:product:layercake}) dan Tonelli untuk
        membuktikan
        \[
        \norm{Mf}_p^p \;\leq\;
        \frac{6p\,2^{p-1}}{p-1}\,\norm f_p^p .
        \]
        (Sedangkan peledakannya saat $p \downarrow 1$ adalah kegagalan pertanyaan
        13, yang \omterm{def:b3:lebesgue:measurable}{terukur}.)
  \item (Teorema pendiferensialan Lebesgue) Buktikan: untuk $f \in
        L^1(\R)$,
        \[
        \frac1{2r}\int_{x-r}^{x+r}f\,\dd\lambda
        \;\longrightarrow\; f(x)
        \qquad (r \to 0)\quad\text{untuk hampir setiap }x .
        \]
        \emph{(Jelas untuk $f$ yang \omterm{def:b3:topology:continuity}{kontinu}. Secara umum tulislah $f
        = g + h$, dengan $g$ \omterm{def:b3:topology:continuity}{kontinu} berpendukung \omterm{def:b3:topology:compact}{kompak} dan
        $\norm h_1 < \varepsilon$
        (\cref{thm:b3:lp:density}); maka himpunan yang
        $\limsup_{r\to0}$ osilasi rata-ratanya melampaui
        $\delta$ terletak di dalam $\{Mh > \delta/2\} \cup \{\abs h
        > \delta/2\}$, yang berukuran $O(\varepsilon/\delta)$;
        lalu biarkan $\varepsilon \to 0$, lalu $\delta \to 0$ sepanjang sebuah
        barisan.)}
  \item Turunkan: (a) hampir setiap titik merupakan \emph{titik
        Lebesgue} bagi $f$; (b) untuk $f \in L^1$, antiturunan
        $F(x) = \int_0^xf$ dapat diturunkan hampir di mana-mana dengan $F' =
        f$ hampir di mana-mana --- yakni separuh integral teorema
        fundamental kalkulus di dunia Lebesgue, yang menutup
        lingkaran yang dibuka tangganya
        (\cref{pb:b3:measure:1}), yang menunjukkan bahwa separuh
        sebaliknya dapat gagal.
  \item (Titik \omterm{ex:b3:lebesgue:gamma}{kepadatan}) Untuk $A \subseteq \R$ yang \omterm{def:b3:lebesgue:measurable}{terukur},
        tunjukkan bahwa hampir setiap $x \in A$ memenuhi
        $\frac{\lambda(A\cap\intcc{x-r}{x+r})}{2r} \to 1$:
        sehingga himpunan \omterm{def:b3:lebesgue:measurable}{terukur} bersifat penuh secara lokal di hampir semua
        titiknya. Lalu sketsakan, dalam dua baris, bagaimana ini melahirkan
        satu lagi bukti teorema Steinhaus
        (\cref{exo:b3:measure:8}).
\end{enumerate}

\medskip
\noindent\textbf{Bagian VI --- Ragam pada tema
perata-rataan.}
\begin{enumerate}[resume]
  \item (Hardy berbobot) Untuk $\alpha < p - 1$ dan $f \geq 0$,
        tunjukkan bahwa
        \[
        \int_0^\infty\Bigl(\frac{F(x)}x\Bigr)^{p}
        x^{\alpha}\,\dd x \;\leq\;
        \Bigl(\frac{p}{p - 1 - \alpha}\Bigr)^{p}
        \int_0^\infty f(x)^p\,x^{\alpha}\,\dd x
        \]
        lewat pengintegralan parsial yang sama, lalu periksalah bahwa
        perbatasan $\alpha = p - 1$ sungguh terlarang
        (sesuaikan contoh penyangkal pertanyaan 10).
  \item (Adjoinnya) Misalkan $H^*f(x) =
        \int_x^{\infty}\frac{f(t)}t\,\dd t$. Tunjukkan bahwa $\langle
        Hf, g\rangle = \langle f, H^*g\rangle$ untuk
        $f, g$ yang tak negatif (lewat Tonelli), lalu buktikan
        $\norm{H^*f}_p \leq p\,\norm f_p$ \emph{(langsung lewat
        pengintegralan parsial, atau dari Hardy pada eksponen sekawannya lewat
        dualitas --- perhatikanlah eksponen mana yang memungut konstanta mana)}.
  \item (Sebuah ketaksamaan bertipe Hilbert) Turunkan bahwa untuk
        $f \in L^p$, $g \in L^q$ yang tak negatif:
        \[
        \int_0^\infty\!\!\int_0^\infty
        \frac{f(x)\,g(y)}{\max(x,y)}\,\dd x\,\dd y
        \;\leq\; (p + q)\,\norm f_p\,\norm g_q
        \]
        \emph{(pisahkanlah sepanjang $y < x$ / $y \geq x$: sebab masing-masing separuhnya
        merupakan pemasangan satu fungsi terhadap transformasi Hardy
        yang lain)}.
  \item (Keoptimalan, secara diskret) Tunjukkan bahwa konstanta
        $\bigl(\frac p{p-1}\bigr)^p$ pada pertanyaan 8 juga
        optimal: ujilah pada $a_n = n^{-1/p}\,\mathbf 1_{n \leq
        N}$, bandingkan kedua ruasnya dengan integral, lalu biarkan $N
        \to \infty$ (yakni cermin diskret Bagian II).
  \item (Sintesis) Tiga operator perata-rataan muncul pada
        soal ini: yakni $H$ milik Hardy, rata-rata Cesàro
        diskret, dan operator maksimal $M$. Nyatakan dalam satu
        baris masing-masing apa yang dikatakan keterbatasannya, amati bahwa ketiganya
        gagal tepat di $p = 1$, lalu terangkan mengapa itu
        kegagalan yang sama tiga kali (yakni ekor harmonik
        $\frac1x$).
\end{enumerate}

\medskip
\noindent\textbf{Bagian VII --- Ketaksamaan Carleman, dan seberapa
tajamkah tajam itu.}
\begin{enumerate}[resume]
  \item (Ketaksamaan Carleman) Misalkan $a_n \geq 0$ dengan $\sum
        a_n < \infty$. Terapkan ketaksamaan Hardy diskret pada
        pertanyaan 8 ke $b_n = a_n^{1/p}$, pakailah
        ketaksamaan rata-rata hitung--ukur, lalu biarkan $p \to
        \infty$ (dengan menunjukkan bahwa $p \mapsto
        \bigl(\frac p{p-1}\bigr)^p$ turun menuju $\eu$) untuk
        memperoleh
        \[
        \sum_{n\geq1}\bigl(a_1a_2\cdots a_n\bigr)^{1/n}
        \;\leq\; \eu\,\sum_{n\geq1}a_n :
        \]
        sehingga rata-rata ukur sebuah barisan terjumlahkan tetap
        terjumlahkan, dengan biaya paling banyak $\eu$.
  \item (Konstanta $\eu$ bersifat optimal) Ujilah $a_n =
        \frac1n\,\mathbf 1_{n\leq N}$: dengan memakai pengapitan
        Stirling \cref{pb:b3:product:1}, tunjukkan bahwa
        $(n!)^{-1/n} = \frac\eu n\bigl(1 +
        O\bigl(\frac{\ln n}n\bigr)\bigr)$, turunkan bahwa kedua
        ruas Carleman tumbuh seperti $\eu\ln N$, lalu simpulkan
        bahwa tak ada konstanta yang lebih kecil daripada $\eu$ yang dapat bekerja. (Amatilah
        polanya: sebab pengoptimal Hardy maupun Carleman
        sama-sama barisan bertipe harmonik yang tepat gagal
        berada di ruangnya.)
  \item (Seberapa lambat ``tajam'' itu didekati?) Ambil $p = 2$.
        Untuk $f = \mathbf 1_{\intcc01}$, hitunglah
        $\norm{Hf}_2/\norm f_2 = \sqrt2$, terhadap batasnya
        $2$. Sedangkan untuk pengoptimal hampiran $f_A$ pada pertanyaan 5,
        buktikan identitas persisnya
        \[
        \norm{Hf_A}_2^2 = 4\ln A - 8 + \frac{8}{\sqrt A},
        \qquad\text{sehingga}\qquad
        \frac{\norm{Hf_A}_2^2}{\norm{f_A}_2^2}
        = 4 - \frac{8 - 8A^{-1/2}}{\ln A} .
        \]
        Nilaikanlah di $A = \eu^{10}$ (dengan nisbah $\approx 1.790$)
        lalu berkomentarlah: sebab supremumnya $2$ didekati dengan laju
        $1/\ln A$ saja --- jadi sebuah konstanta optimal dapat hampir
        tak terlihat secara numerik.

\end{enumerate}
\end{problem}
